\documentclass[12pt]{article}
\usepackage[a4paper,margin=1in]{geometry}
\usepackage[T1]{fontenc}
\usepackage{lmodern,microtype}
\usepackage{amsmath,amssymb,amsthm,mathtools}
\usepackage{booktabs,array,enumitem,xcolor}
\usepackage[numbers,sort&compress]{natbib}
\usepackage[colorlinks=true,linkcolor=blue!55!black,citecolor=blue!55!black,urlcolor=blue!55!black]{hyperref}
\linespread{1.07}

\newtheorem{theorem}{Theorem}[section]
\newtheorem{proposition}[theorem]{Proposition}
\newtheorem{corollary}[theorem]{Corollary}
\theoremstyle{remark}
\newtheorem{remark}[theorem]{Remark}
\newcommand{\cL}{\mathcal L}
\newcommand{\norm}[1]{\left\lVert#1\right\rVert}

\title{Source--Observation Riesz Channels\\
\large Spectral Projectors, Transfer Residues, and Rank-One Physical Modes}
\author{Liang Wang}
\date{July 2026}

\begin{document}
\maketitle

\begin{abstract}
RH-194 finds that every surviving temporal root lies near one genuine base
eigenvalue, but RH-192 forbids interpreting that root as the complete
Frobenius eigenspace.  This paper gives the exact source-relative object.
For an isolated base spectral projector $P$, define
\[
 X_P=PS,
 \qquad Y_P=P^*O^*.
\]
These are the source and observation states cut out by the same Riesz
channel.  Their Frobenius pairing is
\[
 \langle Y_P,X_P\rangle_F=\operatorname{tr}(OPS),
\]
the contour residue of the scalar transfer function
$g(z)=\operatorname{tr}(O(zI-A)^{-1}S)$.

For pairwise disjoint spectral projectors $P_iP_j=\delta_{ij}P_i$, the cross
pairing is diagonal:
\[
 \langle P_i^*O^*,P_jS\rangle_F
 =\delta_{ij}\operatorname{tr}(OP_iS).
\]
When the selected eigenvalues are simple and the residues are nonzero,
residue normalization gives exact biorthogonal right/left eigenstates of
the matrix left-multiplication dynamics.  A 160-case complex nonnormal audit
verifies projector resolution, cross-channel diagonality, transfer partial
fractions, and normalized biorthogonality with zero failures.

The theorem explains the correct rank-one content of RH-194: one
source--observation direction inside each repeated ambient eigenspace.  A
uniform residue lower bound, interval physical projectors, and all-level
channel transport remain open.
\end{abstract}

\section{Why a channel needs both source and observation}

The source-cyclic space removes directions that cannot be reached from $S$.
It may still contain modes invisible to the physical observation.  A
spectral packet intended for determinant or trace data must therefore be
typed by both ends of the experiment.

Let $A\in\mathbb C^{n\times n}$,
$S\in\mathbb C^{n\times m}$, and
$O\in\mathbb C^{m\times n}$.  The scalar transfer function is
\begin{equation}\label{eq:transfer}
 g(z)=\operatorname{tr}\!\left(O(zI-A)^{-1}S\right).
\end{equation}
For large $z$ it has the moment expansion
\begin{equation}\label{eq:moment-expansion}
 g(z)=\sum_{q\ge0}\frac{\operatorname{tr}(OA^qS)}{z^{q+1}}.
\end{equation}
This is precisely the Frobenius source-observation transfer preserved by the
RH-193 cyclic restriction.

\section{Riesz projectors and matrix states}

Let $\Gamma$ be a positively oriented contour in the resolvent set of $A$
and define
\begin{equation}\label{eq:riesz}
 P=\frac{1}{2\pi i}\int_\Gamma(zI-A)^{-1}\,dz.
\end{equation}
Then $P^2=P$ and $PA=AP$.  On the Frobenius space let
$\cL_AX=AX$.

\begin{proposition}[Projected source channel]\label{prop:source-channel}
The state
\begin{equation}\label{eq:right-state}
 X_P=PS
\end{equation}
belongs to the source-cyclic space and satisfies
\begin{equation}\label{eq:right-channel-action}
 \cL_A^qX_P=P A^qS
 \quad(q\ge0).
\end{equation}
The cyclic space generated by $X_P$ is the part of the source orbit whose
base spectrum lies inside $\Gamma$.
\end{proposition}

\begin{proof}
Since $P$ is a polynomial in $A$ in finite dimension, $PS$ belongs to the
span of powers of $A$ applied to $S$.  Commutation gives the orbit identity.
\end{proof}

The dual state is
\begin{equation}\label{eq:left-state}
 Y_P=P^*O^*.
\end{equation}
It is generated by the adjoint dynamics and the observation seed.

\section{The residue pairing}

\begin{theorem}[Channel pairing equals transfer residue]\label{thm:residue}
For every Riesz projector $P$,
\begin{equation}\label{eq:pairing-residue}
 \boxed{
 \langle Y_P,X_P\rangle_F
 =\operatorname{tr}(OPS).
 }
\end{equation}
Moreover
\begin{equation}\label{eq:contour-residue}
 \operatorname{tr}(OPS)
 =\frac{1}{2\pi i}\int_\Gamma g(z)\,dz.
\end{equation}
\end{theorem}

\begin{proof}
Using $P^2=P$,
\[
 \langle P^*O^*,PS\rangle_F
 =\operatorname{tr}\big((P^*O^*)^*PS\big)
 =\operatorname{tr}(OP^2S)
 =\operatorname{tr}(OPS).
\]
Insert the Riesz integral and use linearity and cyclicity of the trace to
obtain \eqref{eq:contour-residue}.
\end{proof}

This identity is invariant under eigenvector scaling and does not require a
choice of basis inside the spectral subspace.

\section{Disjoint channels}

Let $P_1,\ldots,P_r$ be Riesz projectors for pairwise disjoint spectral
sets.  Then $P_iP_j=0$ for $i\ne j$.

\begin{theorem}[Cross-channel diagonality]\label{thm:cross-diagonal}
Put $X_j=P_jS$, $Y_i=P_i^*O^*$, and
$r_i=\operatorname{tr}(OP_iS)$.  Then
\begin{equation}\label{eq:cross-diagonal}
 \boxed{\langle Y_i,X_j\rangle_F=\delta_{ij}r_i.}
\end{equation}
\end{theorem}

\begin{proof}
Directly,
\[
 \langle Y_i,X_j\rangle_F
 =\operatorname{tr}(OP_iP_jS)
 =\delta_{ij}\operatorname{tr}(OP_iS).
\]
\end{proof}

Thus the observation automatically supplies the correct dual annihilation
between distinct spectral channels.  Poor conditioning can remain because
$|r_i|$ may be small relative to $\norm{X_i}\norm{Y_i}$, but there is no
off-diagonal algebraic contamination.

\section{Simple eigenvalues}

Suppose $\lambda$ is simple, with right and left eigenvectors normalized by
$w_\lambda^*v_\lambda=1$.  Then
\begin{equation}\label{eq:rank-one-projector}
 P_\lambda=v_\lambda w_\lambda^*,
\end{equation}
and
\begin{align}
 X_\lambda&=v_\lambda(w_\lambda^*S),\\
 Y_\lambda&=w_\lambda(Ov_\lambda)^*.
\end{align}
Both are single matrix vectors even though the full Frobenius eigenspace has
dimension $m$.

\begin{corollary}[Exact left/right eigenstates]\label{cor:eigenstates}
For a simple eigenvalue,
\begin{equation}\label{eq:eigenstate-actions}
 \cL_AX_\lambda=\lambda X_\lambda,
 \qquad
 \cL_A^*Y_\lambda=\overline\lambda Y_\lambda.
\end{equation}
The scalar residue is
\begin{equation}\label{eq:simple-residue}
 r_\lambda=(w_\lambda^*S)(Ov_\lambda).
\end{equation}
\end{corollary}

The mode is source-observable precisely when $r_\lambda\ne0$.  For a simple
mode this implies both $P_\lambda S\ne0$ and
$P_\lambda^*O^*\ne0$.

\section{Residue-normalized biorthogonality}

Choose simple disjoint eigenvalues $\lambda_1,\ldots,\lambda_r$ with
$r_i\ne0$.  Define columns in the Frobenius state by
\begin{equation}\label{eq:raw-frames}
 V_i=X_i,
 \qquad
 W_i=\frac{Y_i}{\overline{r_i}}.
\end{equation}

\begin{theorem}[Exact source-channel packet]\label{thm:packet}
The frames $V=[V_1,\ldots,V_r]$ and $W=[W_1,\ldots,W_r]$ satisfy
\begin{equation}\label{eq:biorthogonality}
 W^*V=I_r,
\end{equation}
and
\begin{equation}\label{eq:exact-compression}
 \cL_AV=V\Lambda,
 \qquad
 \cL_A^*W=W\Lambda^*,
 \qquad
 W^*\cL_AV=\Lambda,
\end{equation}
where $\Lambda=\operatorname{diag}(\lambda_1,\ldots,\lambda_r)$.
\end{theorem}

\begin{proof}
Equation \eqref{eq:cross-diagonal} gives $W^*V=I$.  The eigenstate actions
hold columnwise, and compression follows.
\end{proof}

The packet has zero right and left residual.  Its rank is the number of
selected source-observable channels, not the full ambient Riesz rank.

\section{Transfer partial fractions}

If $A$ has a complete simple spectrum, summing its rank-one projectors gives
$I$ and
\begin{equation}\label{eq:resolvent-partial-fractions}
 (zI-A)^{-1}=\sum_j\frac{P_j}{z-\lambda_j}.
\end{equation}
Taking the source-observation trace yields
\begin{equation}\label{eq:transfer-partial-fractions}
 \boxed{g(z)=\sum_j\frac{r_j}{z-\lambda_j}.}
\end{equation}
Zero residues are exact pole cancellations in the scalar transfer function.
The minimal source-observable realization keeps precisely the noncanceled
Jordan data; in the simple case it keeps the indices with $r_j\ne0$.

For a selected subset, equation \eqref{eq:transfer-partial-fractions}
splits into the packet contribution and a complementary transfer remainder.
This is a transfer decomposition, not a claim that the full Frobenius
determinant factors with rank $r$.

\section{Complex nonnormal identity audit}

The implementation uses dimensions 2--9, source widths 1--4, and five
independent trials per pair, giving 160 cases.  Each test matrix is similar
to a diagonal matrix with distinct complex eigenvalues and bounded
similarity condition.

The audit checks:
\begin{enumerate}[leftmargin=2em]
\item resolution of the identity by computed projectors;
\item idempotence and cross-projector annihilation;
\item equality of channel pairing and residue;
\item diagonal cross-channel Gram matrix;
\item residue-normalized biorthogonality;
\item equality of the direct transfer solve and modal partial fractions.
\end{enumerate}
All 160 cases pass the $10^{-9}$ aggregate tolerance.

\section{Physical interpretation of RH-194}

RH-194 finds four simple source-observable modes on each side at
$\sigma=0.01$ \citep{WangRH194}.  The present theorem upgrades their
interpretation:
\begin{equation*}
 \text{one temporal root}
 \quad\rightsquigarrow\quad
 (P_\lambda S,\,P_\lambda^*O^*,\,r_\lambda).
\end{equation*}
This triple is basis-free and tied to the physical input-output experiment.
It is the correct finite rank-one channel.

The raw residue normalization can be poorly conditioned.  RH-196 constructs
balanced coordinates on the same exact right/left spectral subspaces and
proves their optimal cross-angle conditioning.

\section{Boundary of the theorem}

The Riesz projectors are exact finite objects when the contours are known.
The physical projectors in RH-194 are presently computed by floating
eigenvectors, not interval contour integration.  The theorem also supplies
no lower bound on $|r_\lambda|$ or its normalized overlap under refinement.

Finally, selecting a spectral quartet by first seeing its eigenvalues is not
yet an intrinsic dynamical construction.  Gate A requires a rule that finds
and transports the channel without importing the desired spectral answer.
Gates B--E and all claims concerning zeta zeros or RH remain untouched.

\bibliographystyle{plainnat}
\bibliography{references}
\end{document}
